\documentclass[11pt,a4paper]{article}
\usepackage[T1]{fontenc}
\usepackage[utf8]{inputenc}
\usepackage{lmodern,microtype}
\usepackage[margin=27mm,headheight=14pt]{geometry}
\usepackage{amsmath,amssymb,amsthm,mathtools}
\usepackage{enumitem,booktabs,longtable,array}
\usepackage{xcolor,hyperref,bookmark,fancyhdr}
\hypersetup{colorlinks=true,linkcolor=blue!45!black,citecolor=blue!45!black,urlcolor=blue!45!black,
 pdftitle={Strong Hahn Measures on Surreal Workspaces},
 pdfauthor={Research manuscript prepared for Vladimir Reshetnikov}}
\usepackage{aliascnt}
\usepackage[capitalise,noabbrev]{cleveref}
\pagestyle{fancy}\fancyhf{}
\fancyhead[L]{\small Strong Hahn measures}
\fancyhead[R]{\small Atomicity and extension}
\fancyfoot[C]{\thepage}
\setlength{\parskip}{4pt}
\setlength{\parindent}{1.2em}
\setlist{itemsep=3pt,topsep=5pt}
\allowdisplaybreaks[2]
\newtheorem{theorem}{Theorem}[section]
\newaliascnt{lemma}{theorem}
\newtheorem{lemma}[lemma]{Lemma}
\aliascntresetthe{lemma}
\crefname{lemma}{Lemma}{Lemmas}
\Crefname{lemma}{Lemma}{Lemmas}
\newaliascnt{proposition}{theorem}
\newtheorem{proposition}[proposition]{Proposition}
\aliascntresetthe{proposition}
\crefname{proposition}{Proposition}{Propositions}
\Crefname{proposition}{Proposition}{Propositions}
\newaliascnt{corollary}{theorem}
\newtheorem{corollary}[corollary]{Corollary}
\aliascntresetthe{corollary}
\crefname{corollary}{Corollary}{Corollaries}
\Crefname{corollary}{Corollary}{Corollaries}
\theoremstyle{definition}
\newaliascnt{definition}{theorem}
\newtheorem{definition}[definition]{Definition}
\aliascntresetthe{definition}
\crefname{definition}{Definition}{Definitions}
\Crefname{definition}{Definition}{Definitions}
\newaliascnt{example}{theorem}
\newtheorem{example}[example]{Example}
\aliascntresetthe{example}
\crefname{example}{Example}{Examples}
\Crefname{example}{Example}{Examples}
\newaliascnt{question}{theorem}
\newtheorem{question}[question]{Question}
\aliascntresetthe{question}
\crefname{question}{Question}{Questions}
\Crefname{question}{Question}{Questions}
\theoremstyle{remark}
\newaliascnt{remark}{theorem}
\newtheorem{remark}[remark]{Remark}
\aliascntresetthe{remark}
\crefname{remark}{Remark}{Remarks}
\Crefname{remark}{Remark}{Remarks}
\newcommand{\R}{\mathbb R}
\newcommand{\C}{\mathbb C}
\newcommand{\Q}{\mathbb Q}
\newcommand{\N}{\mathbb N}
\newcommand{\No}{\mathbf{No}}
\newcommand{\K}{\mathbb K}
\newcommand{\B}{\mathcal B}
\newcommand{\Pset}{\mathcal P}
\newcommand{\supp}{\operatorname{supp}}
\newcommand{\st}{\operatorname{st}}
\newcommand{\val}{\operatorname{v}}
\newcommand{\Fin}{\operatorname{Fin}}
\newcommand{\coef}[2]{[t^{#1}]#2}
\newcommand{\ind}{\mathbf 1}
\newcommand{\sha}{\texttt{4cf691c7d951e037739d32d9f5c387dcce724f3c}}
\newcommand{\sums}{\mathop{\sum}\nolimits^{\mathrm{H}}}
\newcommand{\prods}{\mathop{\prod}\nolimits^{\mathrm{H}}}
\title{\LARGE\bfseries Strong Hahn Measures on Surreal Workspaces\\[5pt]
\Large Atomicity, exact extension criteria,\\and hidden positivity obstructions}
\author{Research manuscript prepared for Vladimir Reshetnikov}
\date{22 September 2026}
\begin{document}
\maketitle
\thispagestyle{empty}
\begin{abstract}
We study a precise choice of infinite additivity for real and complex Hahn
series, and hence for surreal and surcomplex normal forms: the masses of
every disjoint countable family must be strongly Hahn summable. On every
countably separated measurable space, this requirement forces an unexpectedly
rigid representation. Every coefficient measure has finitely many point
atoms, the union of the supports of \emph{all} event masses is automatically
well ordered, and the measure is a strong sum of its singleton masses. No
common-support hypothesis is assumed.

For coherent cylinder data on a countable product of finite sets, we obtain
an exact extension criterion: the combined coefficient support must be well
ordered, and the number of nonzero cylinders at each fixed exponent must
remain bounded with depth. Positivity is a separate issue. We construct
normalized positive cylinder laws whose unique strong signed extension has
a negative singleton mass, and characterize exactly which orders of allowed
supports rule out this phenomenon. For independent finite-state coordinates,
we prove a second exact criterion: only finitely many residue distributions
may be nondeterministic, and the \emph{individual} rare-outcome probabilities
must form a strongly summable family. The resulting law is supported on
finite deviations from a deterministic tail. A three-state example shows
that summability of the total rare-event probabilities is insufficient.

These results are proposed as new statements within the specified summation
framework, not as a first non-Archimedean probability theory or a solution to
a previously named conjecture. Full proofs, limitations, a source audit,
and reproducible finite checks are supplied. Priority remains provisional;
in particular, the interiors of the repository's newly added ZIP archives
were not available for a complete comparison.
\end{abstract}
\vfill
\noindent\textbf{Keywords.} Surreal numbers; surcomplex numbers; Hahn series;
strong summability; vector measures; cylinder measures; Kolmogorov extension;
non-Archimedean probability.
\newpage
{\small\setlength{\parskip}{0pt}\tableofcontents}
\newpage

\section{The problem and the scope of the claims}\label{sec:intro}

\subsection{Why infinite addition is the real issue}
Replacing real probabilities by nonnegative surreal numbers does not specify
how probabilities of countably many disjoint events should be added.
Surreal normal forms provide one particularly exact addition rule: a family
is summable when its monomials have a common admissible support and only
finitely many summands contribute to each monomial. This rule is useful in
formal algebra, but it is much more restrictive than ordinary convergence
of a real series. In particular, even a summable ordinary real series such
as $\sum_{n\ge1}2^{-n}$ is not a strong Hahn sum when its terms are regarded
as constant series.

The repository already emphasizes this distinction and its physical
consequences \cite{Repo,RepoPhysics}. The purpose here is not to repeat the
equal-infinitesimal-weight obstruction. We determine the entire structure
of measures subject to this rule, and then solve the associated extension
problem for arbitrary coherent finite-state cylinder data and for independent
coordinates.

All negative results in this article refer to \emph{strong Hahn
$\sigma$-additivity}, defined in \cref{def:measure}. They do not rule out
finitely additive surreal probabilities, generalized-limit constructions,
coefficientwise real-summed measures, or ultrametrically continuous measures
on other event rings.

\subsection{The questions answered}
The following questions are formulated here as a research problem; they are
not attributed to a named open-problem list.
\begin{question}
Does strong countable additivity itself supply a common well-ordered support
for all event masses, and what atomicity does it impose?
\end{question}
\begin{question}
Which coherent finite-dimensional Hahn-valued laws extend to a strong
measure on the whole Borel $\sigma$-algebra? Is positivity of all finite laws
sufficient for positivity of the extension?
\end{question}
\begin{question}
For independent finite-state coordinates, which infinitesimal probabilities
permit an infinite product law? Can one test only the total probability of
leaving the most likely state?
\end{question}
The answers are, respectively: automatic common support and coefficientwise
finite atomicity; a support-and-width criterion with a genuine additional
positivity obstruction; and an exact individual-rare-entry criterion for
which aggregation is insufficient.

\subsection{Main results and dependencies}
The logical order is
\[
\begin{gathered}
\text{scalar disjoint-finiteness + a support-selection lemma}\\
\Downarrow\\
\text{atomic representation with automatic common support}\\
\Downarrow\\
\text{exact cylinder extension + positivity-order dichotomy},
\end{gathered}
\]
followed by the independent-product criterion, which additionally uses the
classical Neumann support lemma. The atomicity proof does not require
bounded variation, compactness of the sample space, a norm on the Hahn
field, or a preassigned common support.

The principal new-candidate statements are \cref{thm:atomic,thm:cylinder,thm:positivity-order,thm:independent} and the coarse-graining counterexample
in \cref{sec:coarse}. The scalar Boolean-algebra argument, ordinary binary
randomization, Neumann's lemma, and the surreal normal-form representation
are background tools, not novelty claims.

\subsection{Literature and repository comparison}\label{subsec:comparison}
Gonshor's monograph supplies standard surreal normal-form background
\cite{Gonshor}; Neumann's work supplies the ordered-series support machinery
\cite{Neumann}. Neither the existence of Hahn fields nor their usual
summability calculus is claimed here as new.

Benci, Horsten, and Wenmackers develop non-Archimedean probability using,
among other axioms, a homomorphism from functions on finite subsets to a
non-Archimedean field. Their non-Archimedean continuity axiom is different
from requiring finite contribution at every Hahn exponent in every disjoint
countable family \cite[Section 3.1]{BHW}. Consequently, the present atomicity
theorem does not apply to their theory merely because its values are
non-Archimedean.

Ludkovsky and Khrennikov already prove a non-Archimedean analogue of the
Kolmogorov theorem. Their measure definition uses a covering ring,
ultrametric boundedness, and continuity along shrinking families; their
extension theorem concerns a measure-dependent completion of a cylinder
ring \cite[Definitions 2.1--2.4 and Theorem 2.15.2]{LK}. Those are not the
coefficient-local-finiteness axioms on the full Borel $\sigma$-algebra used
here. Thus the phrase ``a first non-Archimedean Kolmogorov theorem'' would be
incorrect for this article.

The repository comparison is pinned to \sha. Its catalog describes
26 maintained reports, while \texttt{docs/new} contains 18 additional ZIP
archives \cite{Repo}. The catalog, repository overview, directory listings,
and a relevant prior physics discussion were inspected. A code search did
not return a Kolmogorov match, but an API search also reported incomplete
indexing, so a negative search is not evidence of exhaustive absence.
The archive interiors, notably \texttt{surreal\_symbolic\_dynamics.zip}, were
not fully retrievable in this session. No theorem in an unread archive is
being declared absent. \Cref{sec:audit} records this boundary explicitly.

\section{Hahn series and the chosen additivity axiom}\label{sec:setup}

\subsection{Conventions}
Throughout, $\N=\{1,2,\ldots\}$; nonnegative word lengths are indicated
explicitly. Let $\Gamma$ be a set-sized linearly ordered abelian group,
written additively, and let $k$ be $\R$ or $\C$. Write
\[
 K=k((t^\Gamma))
 =\left\{\sum_{\gamma\in S}a_\gamma t^\gamma:
       S\subseteq\Gamma\text{ is well ordered},\ a_\gamma\in k\right\}.
\]
Supports are ordered increasingly. For nonzero $f$, put
$\val(f)=\min\supp(f)$ and let $\val(0)=+\infty$. A real Hahn series is
positive exactly when its leading coefficient is positive.

The coefficient field $k$ is embedded as the constant series. When $k=\R$,
the valuation ring of finite elements is
\[
 \mathcal O=\{f:\val(f)\ge0\},\qquad
 \mathfrak m=\{f:\val(f)>0\},\qquad
 \st(f)=\coef{0}{f}\quad(f\in\mathcal O).
\]
Here ``finite'' means bounded in absolute value by an ordinary real number.
No divisibility assumption on $\Gamma$ is needed for the results below.

\begin{definition}[Strong Hahn summability]\label{def:summability}
A set-indexed family $(f_i)_{i\in I}$ in $K$ is \emph{strongly summable} if
\[
 S_I=\bigcup_{i\in I}\supp(f_i)\text{ is well ordered}
 \quad\text{and}\quad
 \#\{i:\coef{\gamma}{f_i}\ne0\}<\infty\quad(\gamma\in\Gamma).
\]
Its sum $\sums_{i\in I}f_i$ is obtained by the finite coefficient sums.
Zero members are permitted, with no restriction on their number.
\end{definition}
Subfamilies and arbitrary regroupings of a strongly summable family remain
strongly summable. For regrouping, the support can only shrink and at a
fixed exponent at most finitely many groups receive a nonzero contribution.
These observations justify all regroupings below once the ungrouped family
has been shown summable.

\begin{lemma}[Positivity of a strong sum]\label{lem:positive-sum}
A strong sum of nonnegative real Hahn series is nonnegative. It is positive
if at least one summand is positive.
\end{lemma}
\begin{proof}
If the nonzero summands are nonempty, their leading exponents have a least
member $\gamma$, because they lie in a well-ordered common support. At
$\gamma$, the only nonzero contributions are leading coefficients of
positive summands. There are finitely many, and their sum is positive.
No smaller exponent occurs.
\end{proof}

\subsection{Measurable spaces}
\begin{definition}
A measurable space $(X,\Sigma)$ is \emph{countably separated} if there are
$E_1,E_2,\ldots\in\Sigma$ such that distinct points of $X$ have different
membership sequences in the $E_n$.
\end{definition}
Every singleton in such a space is measurable: intersect the chosen side
of each $E_n$ containing that point. Every countable subset is consequently
measurable. Standard Borel spaces, including countable products of finite
discrete spaces, have this property. Countable separation does not mean
that $X$ itself is countable.

\begin{definition}[Strong Hahn measure]\label{def:measure}
A \emph{strong Hahn measure} is a map $\mu:\Sigma\to K$ with
$\mu(\varnothing)=0$ such that every pairwise disjoint sequence
$(A_n)_{n\ge1}$ in $\Sigma$ satisfies
\[
 (\mu(A_n))_{n\ge1}\text{ is strongly summable},\qquad
 \mu\left(\bigcup_{n\ge1}A_n\right)=\sums_{n\ge1}\mu(A_n).
 \tag{SH}\label{eq:SH}
\]
For real coefficients it is \emph{positive} if $\mu(A)\ge0$ for all $A$, and
it is a \emph{strong Hahn probability} if also $\mu(X)=1$.
With complex coefficients no order or probability interpretation is imposed.
\end{definition}
Finite additivity follows by adjoining empty sets. We impose neither a
common support for $\mu(\Sigma)$ nor a continuity condition for the intrinsic
valuation topology.

\begin{example}[Three distinct summation phenomena]\label{ex:sums}
For $\Gamma=\Q$, the family $(t^n)_{n\ge1}$ is strong. So is
$(t^{n/(n+1)})_{n\ge1}$, although its finite partial sums do not approach its
Hahn sum in the valuation topology: the remainder always has valuation
less than $1$. In contrast, $(2^{-n}t)_{n\ge1}$ is not strong, because
infinitely many members contribute at exponent $1$. Finally,
$(t^{1/n})_{n\ge1}$ fails the common-support condition.
\end{example}

\section{The coefficient-atomicity mechanism}\label{sec:scalar}

We first isolate an elementary fact about scalar measures. A detailed proof
is useful because no positivity of coefficient measures can be assumed:
a positive Hahn series may have many negative later coefficients.

\begin{lemma}[Infinite Boolean algebras contain disjoint sequences]
\label{lem:boolean}
An infinite Boolean algebra contains a countably infinite family of
pairwise disjoint nonzero elements.
\end{lemma}
\begin{proof}
If it has infinitely many atoms, choose countably many distinct atoms.
Otherwise it has finitely many atoms. If their join is the unit, every
element is a join of a subfamily of these atoms, making the algebra finite.
Thus their complement is nonzero and has no atoms below it. Split this
complement into two nonzero parts, retain one and set the other aside.
The retained part still has no atom below it. Repetition yields the desired
pairwise disjoint nonzero parts. The case with no atoms is the same argument
starting from the unit.
\end{proof}

\begin{lemma}[Disjoint-finite scalar measures are finitely point-atomic]
\label{lem:scalar}
Let $(X,\Sigma)$ be countably separated and let $\lambda:\Sigma\to k$ be
countably additive. Suppose that for every disjoint sequence $(A_n)$,
only finitely many $\lambda(A_n)$ are nonzero. Then
\[
 \lambda(A)=\sum_{x\in F\cap A}c_x\qquad(A\in\Sigma)
 \label{eq:scalar-atomic}\tag{1}
\]
for a finite set $F\subseteq X$ and nonzero $c_x\in k$. The set and
coefficients are uniquely determined after zero coefficients are omitted.
\end{lemma}
\begin{proof}
Call $A\in\Sigma$ \emph{hereditarily null} if $\lambda(B)=0$ for every
measurable $B\subseteq A$. The family $\mathcal N$ of hereditarily null
sets is a $\sigma$-ideal. Indeed, for $B\subseteq\bigcup_n A_n$ with each
$A_n\in\mathcal N$, the disjoint pieces
\[
 B\cap\left(A_n\setminus\bigcup_{j<n}A_j\right)
\]
have measure zero, and countable additivity gives $\lambda(B)=0$.

Consider the quotient Boolean algebra $\Sigma/\mathcal N$. If it were
infinite, \cref{lem:boolean} would supply disjoint nonzero classes.
Choose representatives $C_n$ and replace each by
$C_n\setminus\bigcup_{j<n}C_j$. Their classes remain nonzero, because the
removed intersections are in $\mathcal N$. The resulting representatives
are genuinely disjoint. In each choose a measurable subset $B_n$ with
$\lambda(B_n)\ne0$, contradicting the hypothesis. Therefore the quotient
Boolean algebra is finite.

If the quotient is zero, $\lambda=0$. Otherwise choose disjoint
representatives $A_1,\ldots,A_r$ of its atoms so that their complement is
hereditarily null. For a measurable subset $B\subseteq A_j$, either $B$
or $A_j\setminus B$ is hereditarily null. It follows that
$\lambda(B)$ is either $0$ or $c_j:=\lambda(A_j)$. Moreover $c_j\ne0$,
for otherwise $A_j$ would be hereditarily null.

Let $E_n$ be a countable separating family. For each $j,n$, choose the side
$D_{j,n}$ of $E_n$ for which $A_j\setminus D_{j,n}$ is hereditarily null.
The intersection
\[
 A_j\cap\bigcap_{n\ge1}D_{j,n}
\]
has hereditarily null complement in $A_j$. It cannot be empty because
$c_j\ne0$, and it contains at most one point because the $E_n$ separate
points. Write it as $\{x_j\}$. Thus $\lambda(\{x_j\})=c_j$ and
$A_j\setminus\{x_j\}$ is hereditarily null. This proves (1). Evaluating
singletons proves uniqueness.
\end{proof}

\begin{remark}
This proof uses only a Boolean-algebra observation, countable additivity,
and countable separation. In particular, it does not hide a bounded
variation assumption for a signed or complex coefficient measure.
\end{remark}

\begin{proposition}[Finite coefficient measures]\label{prop:coefficients}
If $\mu$ is a strong Hahn measure on a countably separated space, then for
every $\gamma\in\Gamma$ the scalar set function
\[
 \lambda_\gamma(A)=\coef{\gamma}{\mu(A)}
\]
is a finite linear combination of Dirac measures.
\end{proposition}
\begin{proof}
For a disjoint sequence, (SH) says that the coefficient values are nonzero
at only finitely many indices and that the coefficient of the union is
their finite sum. Thus $\lambda_\gamma$ is countably additive as an ordinary
scalar measure and satisfies \cref{lem:scalar}.
\end{proof}

\section{Automatic common support and complete atomic representation}
\label{sec:atomic}

Finite atomicity at each exponent does not by itself make the union over
all exponents well ordered. The next selection lemma is the missing step.

\begin{lemma}[Simultaneous nonvanishing selection]\label{lem:selection}
Let $D$ be a countable set. For each $n\ge1$, let
\[
 L_n(A)=\sum_{x\in F_n\cap A}c_{n,x}\qquad(A\subseteq D)
\]
where $F_n$ is finite and at least one $c_{n,x}$ is nonzero. There exists
$A\subseteq D$ for which $L_n(A)\ne0$ for infinitely many $n$.
\end{lemma}
\begin{proof}
Choose the membership bits of $A$ independently with ordinary probability
$1/2$. This is only an auxiliary classical randomization: it may be realized
by the binary digits of a uniform number in $[0,1)$ after enumerating $D$.
For each $n$, condition on all bits in $F_n$ except one whose coefficient is
nonzero. The two possible values of $L_n(A)$ differ by that coefficient,
so at least one is nonzero. Therefore
\[
 \mathbb P(L_n(A)\ne0)\ge\tfrac12.
\]
For every $m$, the event $\bigcup_{n\ge m}\{L_n(A)\ne0\}$ has probability
at least $1/2$. Continuity from above of the ordinary real probability
gives
\[
 \mathbb P\bigl(L_n(A)\ne0\text{ infinitely often}\bigr)\ge\tfrac12.
\]
In particular, such an $A$ exists. No Hahn-valued product probability is
used in this argument.
\end{proof}

\begin{theorem}[Atomicity with automatic common support]\label{thm:atomic}
Let $(X,\Sigma)$ be countably separated and let $\mu:\Sigma\to K$ be a
strong Hahn measure. Set $w_x=\mu(\{x\})$. Then:
\begin{enumerate}[label=(\roman*)]
\item The family $(w_x)_{x\in X}$ is strongly summable.
\item For every $A\in\Sigma$,
\[
 \mu(A)=\sums_{x\in A}w_x.
 \label{eq:atomic}\tag{2}
\]
\item The global support
\[
 S(\mu):=\bigcup_{A\in\Sigma}\supp\mu(A)
       =\bigcup_{x\in X}\supp w_x
 \label{eq:global-support}\tag{3}
\]
is well ordered. At every exponent only finitely many $w_x$ have a
nonzero coefficient.
\end{enumerate}
Conversely, every strongly summable family $(w_x)_{x\in X}$ defines by
(2) a strong Hahn measure on \emph{all} subsets of $X$. For real
coefficients this measure is positive if and only if every $w_x\ge0$.
\end{theorem}
\begin{proof}
By \cref{prop:coefficients}, for each $\gamma$ there is a finite set
$F_\gamma$ such that
\[
 \lambda_\gamma(A)=\sum_{x\in A\cap F_\gamma}c_{\gamma,x},
 \qquad c_{\gamma,x}=\coef{\gamma}{w_x}.
 \label{eq:coord-atomic}\tag{4}
\]
Let $S=\{\gamma:\lambda_\gamma\ne0\}$. Suppose $S$ were not well ordered.
In ZFC a linear order which is not well founded contains a strictly
decreasing sequence; choose
$\gamma_1>\gamma_2>\cdots$ in $S$.
The set $D=\bigcup_n F_{\gamma_n}$ is countable. By
\cref{lem:selection}, it has a subset $A$ for which
$\lambda_{\gamma_n}(A)\ne0$ for infinitely many $n$. Since every subset of
$D$ is measurable, $\mu(A)$ is a Hahn series. But its support contains an
infinite strictly decreasing subsequence of the $\gamma_n$, a contradiction.
Thus $S$ is well ordered.

Equation (4) also shows that the nonzero coefficients at any fixed exponent
occur at only finitely many singleton masses. Hence $(w_x)$ is strongly
summable, with common support $S$. Equation (4) now proves (2)
coefficientwise. The two equalities in (3) follow because any nonzero
coefficient measure has a nonzero singleton coefficient.

Conversely, given a strongly summable family, every subfamily is summable.
For a disjoint sequence of subsets, group the terms according to the subset
containing their index. The regrouped family is strong and has the same
sum, proving (SH). Positivity follows from \cref{lem:positive-sum}, and
its necessity follows by evaluating singletons.
\end{proof}

\begin{corollary}[Coefficientwise criterion without a support hypothesis]
\label{cor:coefficient-test}
For an arbitrary map $\mu:\Sigma\to K$ on a countably separated space,
the following are equivalent: $\mu$ is a strong Hahn measure; and each
coefficient function $A\mapsto\coef{\gamma}{\mu(A)}$ is a finite linear
combination of point masses. In the second condition no common support and
no strong additivity are assumed in advance.
\end{corollary}
\begin{proof}
Necessity is \cref{prop:coefficients}. For sufficiency, the proof of
\cref{thm:atomic} from equation (4) onward uses only finite coefficient
atomicity and the existence of the Hahn value $\mu(A)$ for every measurable
$A$. The selection lemma therefore supplies common support without using
(SH). It gives representation (2), whose converse clause proves (SH).
\end{proof}

For an already strong measure, the support of its singleton masses can
also be checked more directly: a descending support sequence would be
witnessed by countably many singleton masses, whose family is strong by
(SH). The selection argument proves the stronger coefficientwise criterion
above, where that shortcut is not available.

\begin{corollary}[Canonical extension and arbitrary disjoint additivity]
\label{cor:all-subsets}
A strong Hahn measure on a countably separated measurable space has a
unique strong extension to $\Pset(X)$. On that extension it is additive
for every set-indexed disjoint family, with the sum interpreted strongly.
\end{corollary}
\begin{proof}
Existence and arbitrary regrouping follow from (2). Any other strong
extension has the same singleton masses, and \cref{thm:atomic}, applied
to $\Pset(X)$, forces the same formula.
\end{proof}
The assertion is stronger than countable additivity, but depends on the
very restrictive summability axiom. It is not an extension theorem for
ordinary Lebesgue measure.

\begin{corollary}[Finite real shadow]\label{cor:shadow}
For a positive strong Hahn probability, every event mass is finite and
$\st\circ\mu$ is an ordinary real probability with finite point support.
\end{corollary}
\begin{proof}
Finite additivity and positivity imply $0\le\mu(A)\le1$. Thus no event
mass has a negative exponent. Its constant coefficient is nonnegative,
and taking constant coefficients gives a scalar probability. By
\cref{prop:coefficients}, that scalar probability has finite point support.
\end{proof}

\begin{corollary}[Cardinality of the atomic set]\label{cor:cardinality}
If $W=\{x:w_x\ne0\}$, then
\[
 W=\bigcup_{\gamma\in S(\mu)}\{x:\coef{\gamma}{w_x}\ne0\}.
\]
Each set in this union is finite. Consequently $W$ is at most countable
when $S(\mu)$ is at most countable, and has cardinal at most
$\max(\aleph_0,|S(\mu)|)$ in general.
\end{corollary}
This does not assert that every strong surreal measure has countably many
atoms: a surreal support can have uncountable order type.

\begin{example}[Countable separation cannot be dropped]\label{ex:cocountable}
Let $X$ be uncountable and let $\Sigma$ be its countable--cocountable
$\sigma$-algebra. Put $\lambda(A)=0$ for countable $A$ and $\lambda(A)=1$
for cocountable $A$. In a disjoint countable family there is at most one
cocountable member; if none occurs the union is countable. Thus $\lambda$
is a strong measure, even with values in the constant real series.
Yet $\lambda(\{x\})=0$ for every $x$, while $\lambda(X)=1$.
This space is not countably separated: the union of the countable exceptional
sides of a countable family of its measurable sets is countable, leaving
uncountably many points with the same membership pattern.
\end{example}

\section{An exact extension theorem for cylinder data}\label{sec:cylinder}

Let $E_n$ be nonempty finite sets and let
\[
 \Omega=\prod_{n\ge1}E_n,\qquad
 T_n=\prod_{j=1}^{n}E_j,\qquad T_0=\{\varnothing\}.
\]
For a word $s\in T_n$, $[s]$ denotes its cylinder. The product Borel
$\sigma$-algebra is countably separated and generated by these cylinders.
Suppose $m_s\in K$ are coherent:
\[
 m_s=\sum_{a\in E_{n+1}}m_{sa}\qquad(s\in T_n).
 \label{eq:coherence}\tag{5}
\]
This condition gives a finitely additive law on the cylinder algebra.
It does not yet define masses of arbitrary Borel events.

For each exponent define its \emph{active width} at depth $n$ by
\[
 r_\gamma(n)=\#\{s\in T_n:\coef{\gamma}{m_s}\ne0\},
 \qquad
 S_T=\bigcup_{n\ge0}\bigcup_{s\in T_n}\supp(m_s).
 \label{eq:width}\tag{6}
\]

\begin{lemma}[Bounded-width coherent scalar trees]\label{lem:tree}
Let $c_s\in k$ satisfy (5). If the number of nonzero $c_s$ at each level
is bounded independently of the level, then there are finitely many distinct
paths $x_1,\ldots,x_r\in\Omega$ and nonzero constants $d_1,\ldots,d_r$
such that
\[
 c_s=\sum_{j:x_j\in[s]}d_j\qquad\text{for every word }s.
 \label{eq:tree-atoms}\tag{7}
\]
This representation is unique.
\end{lemma}
\begin{proof}
Each nonzero parent has at least one nonzero child. Children of different
parents are distinct, so the widths form a nondecreasing sequence of
integers. A bounded such sequence is eventually constant, say from level
$N$ onward. At every subsequent level, each nonzero parent has exactly one
nonzero child, and every zero parent has only zero children. Indeed, either
an extra nonzero child or nonzero children under a zero parent would increase
the width.

The nonzero nodes at level $N$ therefore continue along finitely many
unique paths. Coherence makes the coefficient constant along each path.
These paths and coefficients represent all levels $n\ge N$; summing
descendants gives the earlier levels. For uniqueness, a sufficiently deep
level separates any finite collection of distinct paths, so its cylinder
values recover their individual coefficients.
\end{proof}

\begin{theorem}[Exact strong cylinder-extension criterion]\label{thm:cylinder}
The coherent family $(m_s)$ has a strong Hahn measure extension to
$\B(\Omega)$ if and only if
\begin{align}
 &S_T\text{ is well ordered},\label{eq:C1}\tag{C1}\\
 &\sup_{n\ge0}r_\gamma(n)<\infty\quad\text{for every }\gamma\in\Gamma.
 \label{eq:C2}\tag{C2}
\end{align}
When it exists, the extension is unique, its global support is $S_T$, and
it extends uniquely to all subsets of $\Omega$.
\end{theorem}
\begin{proof}
Suppose first that $\mu$ is an extension. By \cref{thm:atomic}, its global
support is well ordered and each coefficient measure has finitely many
point atoms. At a fixed level, each nonzero coefficient cylinder contains
at least one of those atoms. The cylinders at that level are disjoint, so
their number is bounded by the number of coefficient atoms. This proves
(C1) and (C2). A nonzero finite atomic coefficient measure is nonzero on
some cylinder, by separating its finitely many atoms. Hence the global
support is exactly $S_T$.

Conversely, for each $\gamma\in S_T$, apply \cref{lem:tree} to
$c_s=\coef{\gamma}{m_s}$ and obtain a finite atomic scalar measure
$\lambda_\gamma$ on $\Omega$. Define, for any subset $A\subseteq\Omega$,
\[
 \mu(A)=\sum_{\gamma\in S_T}\lambda_\gamma(A)t^\gamma.
 \label{eq:construct-extension}\tag{8}
\]
Its support is a subset of the well-ordered set $S_T$, so it is a Hahn
series. For a disjoint sequence of subsets, at any fixed $\gamma$ only
finitely many members can meet the finite atomic support of
$\lambda_\gamma$. Thus the family of event masses is strong, and finite
coefficient additivity proves (SH). Formula (7) gives the required cylinder
values.

For uniqueness, any strong extension has finite atomic coefficient
measures by \cref{thm:atomic}. Each is determined by its cylinder values
by the last argument in \cref{lem:tree}. This also excludes additional
exponents not in $S_T$. The all-subsets assertion follows from
\cref{cor:all-subsets}.
\end{proof}

\begin{remark}[What the criterion does and does not bound]
The width bound is allowed to depend on $\gamma$. No uniform bound on the
number of all atoms is required. Likewise the stabilization level in
\cref{lem:tree} may depend on $\gamma$. These two nonuniformities are
essential: they permit infinitely many infinitesimal scales, and are also
responsible for the positivity obstruction below.
\end{remark}

\begin{corollary}[Exact positivity test]\label{cor:positive-test}
Assume (C1)--(C2), $k=\R$, and $m_\varnothing=1$. Reconstruct the coefficient
measures as above and put
\[
 w_x=\sum_{\gamma\in S_T}\lambda_\gamma(\{x\})t^\gamma.
 \label{eq:reconstructed-atoms}\tag{9}
\]
There is a positive strong probability extension if and only if $w_x\ge0$
for every $x\in\Omega$.
\end{corollary}
\begin{proof}
This is the positivity clause of \cref{thm:atomic}, together with uniqueness
in \cref{thm:cylinder}.
\end{proof}

\section{Positive cylinders can force a negative point mass}
\label{sec:positivity}

\subsection{A rank-one example with complete formulas}
Work in $\R((t^\Q))$ and put
\[
 a_n=\frac{n}{n+1},\qquad
 x=(0,0,0,\ldots),\qquad z=(1,1,1,\ldots),
\]
\[
 y_n=(\underbrace{0,\ldots,0}_{n-1},1,0,0,\ldots)
 \quad(n\ge1).
\]
All these paths are distinct. Assign the weights
\begin{align}
 w_{y_n}&=t^{a_n},\nonumber\\
 w_x&=-t,\label{eq:hidden-weights}\tag{10}\\
 w_z&=1-\sums_{n\ge1}t^{a_n}+t,\nonumber
\end{align}
and give all other paths weight zero. The combined support is
\[
 \{0\}\cup\{a_n:n\ge1\}\cup\{1\},
\]
a well-ordered set of order type $\omega+1$. At each $a_n$ there are two
nonzero point coefficients, at exponent $1$ there are two, and at exponent
$0$ there is one. Thus the weights are strongly summable, and their sum
is $1$. They define a normalized strong signed measure $\nu$.

\begin{proposition}[Every finite-dimensional law is positive]
\label{prop:hidden-positive}
For every cylinder $[s]$, $\nu([s])\ge0$. Nevertheless no positive strong
measure has these cylinder values.
\end{proposition}
\begin{proof}
For $N\ge1$, the all-zero cylinder contains $x$ and exactly the $y_n$ with
$n>N$, but not $z$. Hence
\[
 \nu([0^N])=\sums_{n>N}t^{a_n}-t>0,
 \label{eq:zero-cylinder}\tag{11}
\]
because its leading exponent is $a_{N+1}<1$ with coefficient $1$.
A cylinder not containing $x$ contains at most one $y_n$. If it also
contains $z$, its mass has constant coefficient $1$ and is positive.
Otherwise its mass is either one $t^{a_n}$ or zero. The empty-word cylinder
has mass $1$. Therefore all cylinders are nonnegative, and so are all
finite unions of cylinders.

The unique strong signed extension is $\nu$ by
\cref{thm:cylinder}. Since $\nu(\{x\})=-t<0$, a positive strong extension
is impossible.
\end{proof}
Notice that the widths are at most two at every exponent. The failure is
not unbounded width, lack of a common support, or negative finite laws.
Each finite neighborhood of $x$ contains a positive contribution at an
exponent below $1$, but these contributions escape to later and later
paths. The singleton has none of them.

\subsection{An exact order-type criterion for positivity transfer}
The example has the smallest possible infinite support order type which
can exhibit the phenomenon: an increasing sequence followed by a later
exponent. This observation gives a sharp general statement.

\begin{definition}
A well-ordered set $S$ has \emph{finite initial segments} if
$\{\delta\in S:\delta\le\gamma\}$ is finite for every $\gamma\in S$.
For a well-ordered set, this is equivalent to its order type being at most
$\omega$.
\end{definition}
This is an internal condition on the active support. It does not require
$S$ to be unbounded in $\Gamma$. For instance, $\{n/(n+1):n\ge1\}$ has
finite initial segments even though it is bounded in $\Q$.

\begin{theorem}[Exact support-order dichotomy for positivity]
\label{thm:positivity-order}
Let $S\subseteq\Gamma_{\ge0}$ be well ordered with $0\in S$. The following
are equivalent.
\begin{enumerate}[label=(\roman*)]
\item $S$ has finite initial segments.
\item On binary Cantor space, every normalized strong real signed measure
with global support contained in $S$ and nonnegative cylinder masses is
positive on every Borel set.
\item Every normalized nonnegative coherent binary cylinder family with
combined support contained in $S$ satisfying (C2) has a positive strong
extension.
\end{enumerate}
For (i)$\Rightarrow$(ii), binary alphabets may be replaced by arbitrary
finite alphabets.
\end{theorem}
\begin{proof}
Assume (i) and let a strong signed measure $\mu$ have nonnegative cylinder
masses. Suppose some point weight $w_x$ were negative, and let
$\gamma=\val(w_x)$. There are finitely many active exponents $\delta\le\gamma$
and at each there are finitely many coefficient atoms. Their union is a
finite set $F$. Choose a cylinder containing $x$ and none of the other
points of $F$. Its coefficients at all exponents at most $\gamma$ agree
with those of $w_x$. There are no exponents outside $S$. Thus its leading
coefficient is the negative coefficient of $w_x$ at $\gamma$, contradicting
cylinder positivity. Every point weight is nonnegative, and
\cref{thm:atomic} gives (ii).

The equivalence between (ii) and (iii) follows from
\cref{thm:cylinder}: (C1) is automatic for a subset of $S$, and the strong
extension is unique.

Finally suppose (i) fails. Choose $\gamma\in S$ with infinitely many
predecessors. Since $S$ is well ordered, there is a strictly increasing
sequence $0<\delta_1<\delta_2<\cdots<\gamma$ in $S$. Repeat (10) with
$w_{y_n}=t^{\delta_n}$, $w_x=-t^\gamma$, and
\[
 w_z=1-\sums_{n\ge1}t^{\delta_n}+t^\gamma.
\]
Its support lies in $S$, it is normalized and strong, and the proof of
\cref{prop:hidden-positive} applies verbatim because
$\delta_{N+1}<\gamma$. Thus (ii) fails. This proves the equivalence.
\end{proof}

\begin{remark}[Not an order-completeness argument]
We have not passed to an infimum of positive cylinder values. The proof of
positivity in the finite-initial-segment case isolates finitely many
\emph{coefficient atoms}; the counterexample shows exactly why this finite
isolation can fail when an active exponent has infinitely many predecessors.
\end{remark}

\section{Labeled support calculus for infinite products}\label{sec:neumann}

The independent-product theorem needs more than a scalar valuation bound.
It must control both the support and the number of labeled contributions
at a fixed exponent. We state the classical support result in the form
needed here.

\begin{lemma}[Neumann support lemma, classical]\label{lem:neumann}
Let $S\subseteq\Gamma_{>0}$ be well ordered. The additive monoid
\[
 S^*=\{s_1+\cdots+s_m:m\ge0,\ s_j\in S\}
\]
is well ordered. For each $\gamma\in\Gamma$, only finitely many ordered
finite words $(s_1,\ldots,s_m)$ in $S$ have sum $\gamma$.
\end{lemma}
This is the standard positive-support form of Neumann's lemma
\cite{Neumann}. It holds at arbitrary ordered-group rank; the lengths of
words are not bounded by an Archimedean comparison with a least positive
scale. It is also part of the repository's existing Hahn support machinery
\cite{Repo}. We use it as an imported theorem.

\begin{lemma}[Labeled words and subproducts]\label{lem:labeled}
Let $(a_i)_{i\in I}$ be strongly summable, with every nonzero $a_i$ of
positive valuation. Then the family
\[
 \left(a_{i_1}\cdots a_{i_m}\right)_{
         m\ge0,\ (i_1,\ldots,i_m)\in I^m}
 \label{eq:word-family}\tag{12}
\]
is strongly summable, except that zero products may be discarded and the
empty word occurs just once. Any subfamily of these labeled products is
therefore strongly summable, including the products indexed by finite
subsets of $I$.
\end{lemma}
\begin{proof}
Let $S$ be the union of the positive supports of the $a_i$. The support of
every product lies in $S^*$. At a fixed exponent $\gamma$, Neumann's lemma
gives finitely many exponent words $(s_1,\ldots,s_m)$ summing to $\gamma$.
For each letter $s_j$, only finitely many labels $i_j$ can contribute a
nonzero coefficient, by strong summability of $(a_i)$. Thus only finitely
many labeled coefficient words occur at $\gamma$, and in particular only
finitely many nonzero product terms contribute there. This proves both
requirements of strong summability.
\end{proof}

\begin{remark}[Zero labels]
Infinitely many zero members cause no problem in (12): every nonempty
product containing such a label is the zero series, and zero series
contribute at no exponent. The single empty product is $1$.
\end{remark}

\begin{proposition}[Infinite positive-order products]\label{prop:products}
For a strongly summable positive-order family $(q_n)_{n\ge1}$, define
\[
 U=\prods_{n\ge1}(1-q_n)
   :=\sums_{F\in\Fin(\N_{>0})}(-1)^{|F|}\prod_{n\in F}q_n.
 \label{eq:U}\tag{13}
\]
Then $U=1+\mathfrak m$ and is a unit. The families
\[
 \left(\frac{q_n}{1-q_n}\right)_n
 \quad\text{and more generally}\quad
 \left(\frac{a_{n,j}}{1-q_n}\right)_{n,j}
 \label{eq:ratios}\tag{14}
\]
are strongly summable whenever the positive-order family $(a_{n,j})_{n,j}$
is strong and $q_n=\sum_j a_{n,j}$ with finitely many $j$ at each $n$.
All the finite-factor product identities used with (13) remain valid
coefficientwise.
\end{proposition}
\begin{proof}
Equation (13) is a subfamily of the labeled words of
\cref{lem:labeled}, with signs. Its nonconstant support is positive, hence
$U=1+\mathfrak m$. The geometric series
$(1-q_n)^{-1}=\sums_{m\ge0}q_n^m$ is valid by the same lemma. Expanding
$a_{n,j}\sum_mq_n^m$ gives labeled positive-order words, with the first
letter carrying the distinguished index $(n,j)$. At a fixed exponent
there are only finitely many such words and such distinguished labels.
This proves (14). Every claimed product identity can now be checked at a
fixed exponent, where the relevant words and labels form a finite set;
the identity reduces to the corresponding identity for finitely many
factors.
\end{proof}

\begin{lemma}[Recovering rare entries from first-error masses]
\label{lem:recover}
Suppose $p_{n,j}$ are positive-order real Hahn series, with finitely many
indices $j$ at each $n$, and put
\[
 q_n=\sum_jp_{n,j},\quad u_n=1-q_n,\quad
 h_{n,j}=p_{n,j}\prod_{k<n}u_k.
 \label{eq:first-error}\tag{15}
\]
Assume $q_n\in\mathfrak m$ for all $n$. Then the family $(p_{n,j})_{n,j}$
is strongly summable if and only if $(h_{n,j})_{n,j}$ is strongly summable.
Zero entries are allowed.
\end{lemma}
\begin{proof}
If $(p_{n,j})$ is strong, expansion of the finite products in (15) gives
labeled words whose first letter has the distinguished label $(n,j)$.
\Cref{lem:labeled} shows that the whole family $(h_{n,j})$ is strong.

For the converse, finite telescoping yields
\[
 \prod_{k<n}u_k=1-H_{<n},\qquad
 H_{<n}=\sum_{k<n}\sum_jh_{k,j}.
 \label{eq:survival}\tag{16}
\]
Each nonzero $h_{n,j}$ has positive valuation. Therefore
\[
 p_{n,j}=h_{n,j}(1-H_{<n})^{-1}
        =\sums_{m\ge0}h_{n,j}H_{<n}^{\,m}.
 \label{eq:recover}\tag{17}
\]
Expand $H_{<n}^{\,m}$ into its labeled words. The first label $(n,j)$ is
again distinguished; all subsequent coordinate labels are smaller than
$n$. At each exponent, \cref{lem:labeled} leaves only finitely many
possible full words, hence only finitely many distinguished indices.
Thus (17), simultaneously for all $(n,j)$, is strongly summable. This
proves the converse without imposing a uniform bound on $n$ or on $m$.
\end{proof}

\section{Independent coordinates: a sharp extension criterion}
\label{sec:independent}

Let $p_{n,a}\in\R((t^\Gamma))$ satisfy
\[
 p_{n,a}\ge0,\qquad \sum_{a\in E_n}p_{n,a}=1.
 \label{eq:prob-rows}\tag{18}
\]
In particular, every entry is finite. The independent cylinder law is
\[
 m_s=\prod_{n=1}^{|s|}p_{n,s_n}.
 \label{eq:independent-cylinder}\tag{19}
\]
Write $\overline p_{n,a}=\st(p_{n,a})$. This is an ordinary finite-state
probability vector. Call coordinate $n$ \emph{residually nondeterministic}
if at least two entries of this vector are positive. Let $J$ be the set
of such coordinates. If $n\notin J$, there is a unique $b_n\in E_n$ with
$\overline p_{n,b_n}=1$. Every other entry is zero or a positive
infinitesimal; call these the \emph{rare entries}.

\begin{theorem}[Exact independent finite-state criterion]
\label{thm:independent}
The independent cylinder law (19) has a positive strong Hahn probability
extension if and only if
\begin{enumerate}[label=(\roman*)]
\item $J$ is finite;
\item the indexed family
\[
 \bigl(p_{n,a}\bigr)_{n\notin J,\ a\ne b_n}
 \label{eq:rare-family}\tag{20}
\]
is strongly summable.
\end{enumerate}
When these conditions hold, the extension is unique and is supported on
paths with only finitely many deviations from $b_n$ outside $J$.
\end{theorem}

\subsection{Necessity: finite residue branching}
\begin{proof}[Proof of necessity in \cref{thm:independent}]
Assume a positive strong extension $\mu$ exists. By \cref{cor:shadow},
$\st\circ\mu$ is a real probability with a finite set of point atoms, say
at most $M$. If there were $r$ residually nondeterministic coordinates,
choose two positive-residue states at each of them. The corresponding
$2^r$ disjoint finite-coordinate cylinder events all have positive
constant coefficients by (19). Each must contain a distinct residue atom.
Thus $2^r\le M$. It follows that $J$ is finite.

Project away the coordinates in $J$ and relabel the remaining coordinates
in increasing order. Pushforward preserves (SH), so we may assume $J$ is
empty. Put
\[
 q_n=\sum_{a\ne b_n}p_{n,a},\qquad u_n=p_{n,b_n}=1-q_n.
\]
Each $q_n$ is zero or a positive infinitesimal. For $a\ne b_n$, let
$B_{n,a}$ be the event that coordinates $1,\ldots,n-1$ equal their base
states and coordinate $n$ equals $a$. These events are pairwise disjoint
as $(n,a)$ varies, and their prescribed masses are
\[
 \mu(B_{n,a})=p_{n,a}\prod_{k<n}u_k=h_{n,a}.
\]
The family is countable because every alphabet is finite. Strong
countable additivity makes $(h_{n,a})$ strongly summable.
\Cref{lem:recover} therefore makes the individual rare-entry family
$(p_{n,a})_{a\ne b_n}$ strongly summable, proving (ii).
\end{proof}

\subsection{Sufficiency: the finite-deviation formula}
For the construction, reorder the complement of the finite set $J$ as
$1,2,\ldots$ temporarily; the notation below continues to use its original
indices. Define
\[
 q_n=\sum_{a\ne b_n}p_{n,a},\quad
 u_n=1-q_n,\quad r_{n,a}=\frac{p_{n,a}}{u_n}
 \quad(n\notin J,\ a\ne b_n),
\]
and
\[
 U=\prods_{n\notin J}u_n.
 \label{eq:baseproduct}\tag{21}
\]
The $q_n$ form a strong family by grouping (20), and
\cref{prop:products} applies. Let $\Omega_{\mathrm{fin}}$ be the paths
$x$ for which
\[
 F(x):=\{n\notin J:x_n\ne b_n\}
\]
is finite. Put
\[
 W_x=
 \left(\prod_{j\in J}p_{j,x_j}\right)
 U\prod_{n\in F(x)}r_{n,x_n}
 \quad(x\in\Omega_{\mathrm{fin}}),
 \qquad W_x=0\quad(x\notin\Omega_{\mathrm{fin}}).
 \label{eq:product-weights}\tag{22}
\]
The product over $J$ is ordinary and finite. Thus residue randomness is
retained exactly at the finitely many exceptional coordinates; it is not
replaced by a chosen state there.

\begin{proof}[Proof of sufficiency and uniqueness in \cref{thm:independent}]
The family $(r_{n,a})$ is strong and positive-order by
\cref{prop:products}. Its products indexed by finite coordinate subsets
with one rare-state label per chosen coordinate are a subfamily of the
labeled word family. They are strongly summable. Multiplication by the
fixed series $U$ preserves strong summability: the sum of its support with
the common support is well ordered and has finite coefficient
contributions, by the usual Hahn multiplication rule. Taking the finite
union corresponding to the choices at coordinates in $J$ and multiplying
by their fixed finite products preserves this property. Consequently
$(W_x)$ is strong.

Every $u_n$ has leading coefficient $1$, so $U$ also has leading
coefficient $1$ and is positive. Each other factor in (22) is nonnegative.
Thus every $W_x\ge0$.

For $n\notin J$,
\[
 1+\sum_{a\ne b_n}r_{n,a}
 =1+\frac{q_n}{u_n}=u_n^{-1}.
 \label{eq:one-plus-r}\tag{23}
\]
Summing the finite-deviation products and applying the justified
coefficientwise finite-factor identities gives
\[
 \sums_{\substack{F\subseteq\N\setminus J\text{ finite}\\
                         a_n\ne b_n\ (n\in F)}}
       \prod_{n\in F}r_{n,a_n}
 =\prods_{n\notin J}\left(1+\sum_{a\ne b_n}r_{n,a}\right)=U^{-1}.
 \label{eq:normalization-product}\tag{24}
\]
The finite sum over the coordinates in $J$ of
$\prod_{j\in J}p_{j,x_j}$ is $1$. Equations (22)--(24) show
$\sums_x W_x=1$.

By \cref{thm:atomic}, these weights define a positive strong probability
on all subsets. To check a finite-coordinate cylinder, separate the
specified coordinates from the unspecified ones in (24). At a specified
base coordinate the factor is $u_n$; at a specified rare coordinate it is
$u_nr_{n,a}=p_{n,a}$. At an unspecified coordinate the summed contribution
is $u_n(1+\sum_a r_{n,a})=1$. The finite exceptional coordinates satisfy
the same marginal identity by (18). The resulting cylinder mass is
exactly (19).

Uniqueness follows from \cref{thm:cylinder}, or directly from uniqueness
of the finite atomic coefficient measures determined by cylinders. Formula
(22) proves concentration on $\Omega_{\mathrm{fin}}$.
\end{proof}

\begin{remark}[A precise meaning of concentration]
The set $\Omega_{\mathrm{fin}}$ is countable, because the alphabets and $J$
are finite and only a finite set of other coordinates may vary. It is
Borel. Formula (22) gives
$\mu(\Omega\setminus\Omega_{\mathrm{fin}})=0$ and
$\mu(\Omega_{\mathrm{fin}})=1$. An infinite-deviation path can still lie
in every one of its positive-mass finite cylinders. Such a path need not
be a point atom.
\end{remark}

\subsection{Bernoulli variables}
\begin{corollary}[One-line Bernoulli criterion]\label{cor:bernoulli}
Let $0\le p_n\le1$ be real Hahn series, prescribing independent Bernoulli
coordinates with success probabilities $p_n$. Set
\[
 q_n=\min(p_n,1-p_n).
\]
There is a positive strong Borel product law if and only if the family
$(q_n)_{n\ge1}$ is strongly summable.
\end{corollary}
\begin{proof}
A residue Bernoulli distribution is nondeterministic exactly when
$\st(q_n)>0$. Since every $q_n\ge0$, coefficient-local finiteness at
exponent $0$ implies that this occurs at only finitely many $n$ whenever
$(q_n)$ is strong. Outside those indices, $q_n$ is exactly the rare-entry
probability. Conversely, a strong rare-entry tail together with finitely
many exceptional $q_n$ is a strong family. Apply
\cref{thm:independent}.
\end{proof}
Choose at every coordinate a state of probability $1-q_n$ as the base
state, resolving ties arbitrarily. With the finite noninfinitesimal factors
separated before forming an infinite product, the probability of the path
whose deviations occur exactly at a finite set $F$ is
\[
 \left(\prods_{n\ge1}(1-q_n)\right)
       \prod_{n\in F}\frac{q_n}{1-q_n}.
 \label{eq:Bernoulli-atom}\tag{25}
\]
The denominator never vanishes since $q_n\le1/2$.

\begin{example}[Successes and failures of extension]\label{ex:bernoulli}
The probabilities $p_n=t^n$ admit a strong product law; its all-zero
atom is $\prods_{n\ge1}(1-t^n)$. The probabilities
$p_n=t^{n/(n+1)}$ also admit one, even though the associated partial sums
need not converge in the valuation topology.

Neither $p_n=2^{-n}t$ nor $p_n=t$ admits a strong product law: infinitely
many rare entries contribute at exponent $1$. The failure persists despite
ordinary summability of the real coefficients in the first example.
The probabilities $p_n=t^{1/n}$ fail because their supports descend.
Finally $p_n=1/2$ at infinitely many coordinates fails already at the
residue level. Finitely many ordinary nontrivial probabilities together
with a strong infinitesimal rare tail do admit an extension.
\end{example}

\section{Coarse-graining can hide a genuine obstruction}\label{sec:coarse}

The finite-state theorem tests the individual rare entries, not merely
their positive sum at each coordinate. This distinction is necessary
because positive Hahn series can have negative higher coefficients.

Take
\[
 \Gamma=\Q\oplus_{\mathrm{lex}}\Q,\qquad
 \alpha=(0,1),\quad\beta=(1,0),\quad
 \gamma_n=\beta+\frac1n\alpha.
\]
Thus $\gamma_1>\gamma_2>\cdots$, but every $n\alpha<\gamma_n$.
For $E_n=\{0,1,2\}$ put
\begin{align}
 p_{n,0}&=1-t^{n\alpha},\nonumber\\
 p_{n,1}&=\tfrac12\bigl(t^{n\alpha}+t^{\gamma_n}\bigr),
 \label{eq:three-state}\tag{26}\\
 p_{n,2}&=\tfrac12\bigl(t^{n\alpha}-t^{\gamma_n}\bigr).\nonumber
\end{align}
All three entries are positive and their sum is $1$. The residue state is
$0$ with probability $1$.

\begin{theorem}[Extension may fail after refining a valid binary law]
\label{thm:coarse}
The independent three-state cylinder laws (26) have no positive strong
Borel extension. Their binary coarse-graining which merges states $1$ and
$2$ does have a positive strong Borel extension.
\end{theorem}
\begin{proof}
The coarse-grained rare probability is
\[
 p_{n,1}+p_{n,2}=t^{n\alpha}.
\]
Its family is strongly summable, so \cref{cor:bernoulli} supplies a binary
strong product law. The individual rare-entry supports, however, contain
the strictly decreasing sequence $(\gamma_n)$. Thus (20) is not strongly
summable, and \cref{thm:independent} rules out a three-state extension.

There is also a direct witness. Let $B_n$ be the event that the first
$n-1$ coordinates are $0$ and the $n$th coordinate is $1$. These events
are pairwise disjoint and have prescribed masses
\[
 h_n=\tfrac12\bigl(t^{n\alpha}+t^{\gamma_n}\bigr)
                 \prod_{k<n}(1-t^{k\alpha}).
 \label{eq:direct-coarse}\tag{27}
\]
The coefficient of $t^{\gamma_n}$ in $h_n$ is $1/2$. Contributions from
$t^{n\alpha}$ have first coordinate $0$, and the terms from
$t^{\gamma_n}$ multiplied by a nonempty subproduct have exponent strictly
larger than $\gamma_n$. Therefore the combined support of the $h_n$
contains the descending sequence $(\gamma_n)$, so this disjoint family
cannot satisfy (SH).
\end{proof}

\begin{remark}[Why positivity does not protect support]
The negative coefficient of $t^{\gamma_n}$ in $p_{n,2}$ does not make that
probability negative: the smaller exponent $n\alpha$ has positive leading
coefficient. Upon aggregation, the two higher coefficients cancel. This
is exactly the information lost by testing only total rare probabilities.
The example does not contradict preservation of strong measures by
pushforward. It shows that a valid coarse law need not admit a prescribed
strong refinement.
\end{remark}

\section{Transfer to actual surreal and surcomplex numbers}\label{sec:surreal}

\subsection{The normal-form embedding and its sign convention}
Assume an order-preserving additive embedding $\iota:\Gamma\to\No$ is
fixed. The normal-form map is
\[
 \sum_{\gamma\in S}a_\gamma t^\gamma
       \longmapsto
 \sum_{\gamma\in S}a_\gamma\omega^{-\iota(\gamma)}.
 \label{eq:surreal-map}\tag{28}
\]
The minus sign is essential: increasing well-ordered Hahn exponents become
decreasing well-based surreal exponents. The standard normal-form theorem
identifies (28) with a field embedding for real coefficients; applying it
to real and imaginary parts gives a map into $\No[i]$
\cite{Gonshor}. It preserves the designated strong sums, not an arbitrary
notion of topological convergence.

For $\Gamma=\Q$, take $\iota(\gamma)=\gamma$. Thus the hidden-negative-atom
example is literally a surreal example with $t=\omega^{-1}$.
For the rank-two example, use
\[
 \iota(a,b)=a\omega+b.
\]
This is additive and strictly order preserving. Its monomials are
\[
 t^{n\alpha}\mapsto\omega^{-n},\qquad
 t^{\gamma_n}\mapsto\omega^{-\omega-1/n}.
 \label{eq:surreal-coarse}\tag{29}
\]
Consequently (26) is an explicit system of positive surreal probabilities,
not merely a formal example in an unrelated coefficient field.

\subsection{Localization of a set-sized family of surreal values}
Let $X$ and $\Sigma$ be sets, and let $\mu:\Sigma\to\No$ or $\No[i]$ be
a class function restricted to that set domain. Work in a class theory
with the usual replacement principle, or equivalently in a sufficiently
large universe for the set-sized data. The range is set-sized. Each of
its normal forms has set-sized support, so the union of the supports of
all values is a set, whether or not it is well based.

Take the additive subgroup of $\No$ generated by the negatives of all
these exponents. It is a set-sized ordered abelian group $\Gamma$. Every
value of $\mu$ is then represented in $\R((t^\Gamma))$, or in
$\C((t^\Gamma))$ after combining the two coordinate supports. The
strong-additivity axiom is unchanged. Thus \cref{thm:atomic} applies and
proves a \emph{common well-ordered support} as a theorem, rather than
assuming it during localization.

\begin{corollary}[Actual surreal atomicity]\label{cor:actual}
On a countably separated set-sized measurable space, every surreal-valued
or surcomplex-valued measure satisfying countable additivity by canonical
strong normal-form sums is a strong sum of its singleton masses. Every
normal-form coefficient is carried by finitely many points, and all event
masses have a common admissible support. A positive normalized surreal
measure of this kind has a finitely supported ordinary real standard part.
\end{corollary}
No topology on the proper class $\No$ is required here. The sample space
is a set; the theorem makes no unqualified assertion about a measure on
the full power class of $\No$.

\subsection{Enlarging the workspace cannot repair these obstructions}
\begin{proposition}[Absoluteness under Hahn exponent enlargement]
\label{prop:enlarge}
Let $j:\Gamma\to\Delta$ be an order-preserving additive embedding, and
transport coherent cylinder values from $k((t^\Gamma))$ to
$k((u^\Delta))$ by $t^\gamma\mapsto u^{j(\gamma)}$. A strong extension
exists in the larger field if and only if it exists in the original field.
For real coefficients, the same is true for positive strong extensions.
\end{proposition}
\begin{proof}
The combined support is carried order-isomorphically to its image, so
(C1) is unchanged. The widths at image exponents are unchanged, and the
widths at all other exponents are zero. Therefore (C2) is unchanged.
\Cref{thm:cylinder} proves the equivalence for signed or complex
extensions. Its uniqueness proof shows that an extension in the larger
field cannot acquire extra exponents: a finite atomic coefficient measure
which vanishes on every cylinder is zero. The reconstructed singleton
weights are exactly the transported original weights. Their signs are
preserved by the ordered embedding, proving the positive assertion.
\end{proof}
In particular, allowing more surreal exponents does not rescue any of the
failed cylinder extensions above. A proposed surreal extension has
set-sized range and can be localized together with the original data in
a common larger exponent group; the same argument then applies.

\subsection{Uncountably many infinitesimal atoms are possible}
Atomicity is not finiteness of the entire surreal probability law. To see
this, let $X=\{x_\ast\}\cup\{x_\xi:\xi<\kappa\}$ for any set ordinal
$\kappa$, with all the points distinct. Choose the increasing positive
surreal exponents $\xi+1$ and set
\[
 w_{x_\xi}=\omega^{-(\xi+1)},\qquad
 w_{x_\ast}=1-\sum_{\xi<\kappa}\omega^{-(\xi+1)}.
 \label{eq:many-atoms}\tag{30}
\]
The latter sum is a legitimate infinitesimal normal form. The whole family
is strong, every weight is positive, and the weights sum to $1$.
It therefore defines a positive strong measure on all subsets of $X$.
Its real shadow is the single atom at $x_\ast$, while it has $\kappa$ further
positive infinitesimal atoms. For a countably separated instance with
uncountably many atoms, apply this construction to a well ordering of the
underlying set $\R$ and restrict the resulting measure to its Borel sets. This illustrates the distinction between
finite coefficient atomicity and global finite atomicity.

\section{Integration, functoriality, and boundaries of the theory}
\label{sec:integration}

\begin{proposition}[Pushforward and scalar-valued integration]
\label{prop:integration}
Let $\mu$ have the representation (2).
\begin{enumerate}[label=(\roman*)]
\item Any measurable map $f:X\to Y$ has a strong pushforward measure
$f_*\mu(B)=\mu(f^{-1}(B))$.
\item For every function $g:X\to k$, with no boundedness or measurability
assumption needed after extension to all subsets, the expression
\[
 I_\mu(g)=\sums_{x\in X}g(x)w_x
 \label{eq:integral}\tag{31}
\]
is a well-defined Hahn series. It is $k$-linear, agrees with $\mu$ on
indicators, and satisfies
\[
 \coef{\gamma}{I_\mu(g)}
   =\sum_{x\in F_\gamma}g(x)\coef{\gamma}{w_x}.
\]
\item If $k=\R$, $\mu$ is positive, and $g\ge0$, then $I_\mu(g)\ge0$.
\end{enumerate}
\end{proposition}
\begin{proof}
Inverse images preserve disjoint unions, proving (i) directly from (SH).
Alternatively, group the atomic weights over the fibers of $f$.
Multiplication by $g(x)\in k$ introduces no new exponent. At any exponent
only the finitely many original coefficient atoms can contribute, so the
family in (31) is strong. Its coefficient formula proves linearity and the
indicator identity. Positivity follows from \cref{lem:positive-sum}.
\end{proof}

This is a useful integration operation, but it is not classical Lebesgue
integration with coefficients silently promoted to surreal scalars. Its
apparent integrability of every ordinary scalar function is a consequence
of finite evaluation at each exponent.

\begin{example}[Failure for Hahn-valued integrands]
On $X=\{0,1,2,\ldots\}$ put $w_n=t^n$ for $n\ge1$ and
$w_0=1-\sums_{n\ge1}t^n$. This is a positive strong probability.
The Hahn-valued function $g(n)=t^{-n}$ for $n\ge1$, with $g(0)=0$,
has $g(n)w_n=1$ for every $n\ge1$. The integral family is not strongly
summable. Thus (31) deliberately assumes that $g$ takes values in the
ordinary coefficient field, not arbitrary values in $K$.
\end{example}

\begin{example}[An ordinary measure that is not strong]
Lebesgue probability on $[0,1]$, embedded as constant Hahn series, fails
(SH): choose disjoint intervals of positive lengths $2^{-n}$, with harmless
endpoint conventions. Every mass contributes to exponent $0$.
The ordinary series of lengths converges to $1$, but it is not a strong
Hahn family. Hence the finite-shadow conclusion does not exclude an
ordinary real probability theory embedded by a different addition rule;
it excludes treating its usual infinite sums as strong sums.
\end{example}

A more permissive theory could allow absolutely convergent real coefficient
sums at each exponent, or use topological or generalized-limit
additivity. Such theories have different extension and positivity problems.
Theorems in this article should not be transferred to them without checking
the changed hypotheses. Conversely, replacing (SH) by a weaker rule to
recover diffuse probability would abandon the exact finite-contribution
property proved here to cause the rigidity.

\section{Verification, novelty audit, and further questions}\label{sec:audit}

\subsection{What has been proved and what has been checked}
The manuscript supplies mathematical proofs of the atomic representation,
the cylinder criterion, the positivity-order dichotomy, the independent
criterion, and the coarse-graining obstruction. It imports the standard
Neumann lemma and the standard surreal normal-form field representation.
These imports are separately identified in the text. None of the proofs
is claimed to be formally checked in Lean.

The accompanying \path{code/verify.py} uses exact rational arithmetic and
finite-support series. It checks finite tree coherence and width
monotonicity, the finite independent normalization and marginal identities,
finite versions of the hidden-negative-atom construction, and the exact
coefficient witness (27) in a lexicographic rank-two group. Its recorded
output is in \path{data/verification.txt}; the supplied run passed
5,135 exact assertions. The script is not a proof of
an infinite theorem, cannot establish that a support is well ordered from
a finite prefix, and cannot certify eventual boundedness of infinitely
many widths. These limitations matter especially here.

\subsection{Scope of the novelty search}
The public literature comparison was performed on 22 September 2026.
Primary sources examined include the Benci--Horsten--Wenmackers axioms and
the Ludkovsky--Khrennikov definition and extension theorem, with page images
checked for the relevant definitions. Searches targeted Hahn-valued
probability, strong Hahn additivity, and non-Archimedean Kolmogorov
extension. Some broad searches returned irrelevant results; they provide
no evidence of absence.

The repository audit is a snapshot audit, not an assertion about future
commits. The inspected catalog covers the maintained reports and explicitly
distinguishes strong summability from convergence. A prior physics text
already observes the failure of countably many equal infinitesimal weights.
That observation is not one of the new claims. The 18 files in
\texttt{docs/new} were inventoried by name, but their ZIP interiors were
not all accessible. In particular, titles alone cannot exclude overlap
with \texttt{surreal\_symbolic\_dynamics.zip}, either infinite-spectral
archive, or another newly added manuscript.

\begin{center}
\begin{tabular}{p{0.29\linewidth}p{0.63\linewidth}}
\toprule
Status & Meaning in this manuscript\\
\midrule
Proved here & A complete mathematical argument is included, subject to the
explicit imported background theorems.\\
New candidate & The exact statement was not located in the inspected
sources; no claim of established priority or exhaustive archive exclusion
is made.\\
Classical background & Neumann support calculus, scalar and Boolean
arguments, and surreal normal forms are not claimed as inventions.\\
Finite check & Reproducible exact calculations support formulas and detect
some mistakes, but do not prove the infinite statements.\\
Not claimed & A solution of a famous named open conjecture, complete
repository coverage, a replacement for all probability theories, or a Lean
formalization.\\
\bottomrule
\end{tabular}
\end{center}

The strongest priority candidates are the combination of automatic common
support with atomicity, the exact finite-width extension criterion, the
support-order characterization of positivity transfer, and the individual
rare-entry criterion with the refinement counterexample. An independent
literature review and complete archive comparison could change the novelty
assessment without changing the validity of the proofs.

\subsection{A possible formalization path}
The elementary scalar lemma can be formalized using the $\sigma$-ideal of
hereditarily null sets and its Boolean quotient, or by reusing finite atomic
scalar measure results. The simultaneous-selection lemma can use the
ordinary binary product probability space; it must not be instantiated with
the Hahn measure under construction. The coefficient tree lemma needs only
finite branching, monotonicity of natural-number widths, and stabilization
of a bounded monotone sequence.

The Hahn stages fit an interface providing coefficient maps, strong-family
regrouping, positive-support word finiteness, and geometric inversion.
A suitable Lean development should keep the following data separate:
coherence of cylinders, well ordering of their combined support, boundedness
of each coefficient width, and signs of reconstructed point weights.
Collapsing the last property into cylinder positivity would encode the
false assertion refuted by (10).

\subsection{Further questions, not used in any proof}
One natural continuation is to replace finite alphabets by standard Borel
coordinate spaces. Atomicity still holds, but a finite-width cylinder tree
no longer directly describes the extension problem; a coefficientwise
finite atomic compatibility criterion would need to address how point
locations are recovered from the chosen cylinder family.

A second continuation is a weaker addition rule that permits convergent
ordinary real coefficient sums. It would retain diffuse real shadows but
lose the finite Boolean quotient used in \cref{lem:scalar}. The correct
uniform support and variation conditions for an exact extension theorem
in that framework are not established here.

A third question is an effective version. The present criteria are
necessary and sufficient mathematical conditions, not decision procedures
on arbitrary coefficient oracles. Finite-prefix verification cannot decide
all support well-ordering or eventual-width properties. A useful algorithm
would require a restricted representation with certificates for these
infinite conditions.

\section{Conclusion}
Strong Hahn summation is not merely a convenient way to write an infinite
probability sum. When imposed on all disjoint countable event families, it
forces a complete atomic description: finitely many point coefficients at
each exponent, and an automatically well-ordered common support. This
rigidity leads to an exact cylinder extension theorem and an explicit
infinite product formula, but also to two boundaries invisible at a fixed
finite stage. Positive cylinders need not give a positive strong extension,
and summable aggregate rare probabilities need not give summable individual
rare entries.

The results apply to actual surreal and surcomplex values through their
normal forms. Their force comes from the summation axiom, not from any
claim that non-Archimedean probability in general is impossible. They offer
a precise answer to what can, and what cannot, be built when strong
normal-form summation is the intended infinite addition.

{\small
\begin{thebibliography}{9}
\setlength{\itemsep}{3pt}
\bibitem{Gonshor}
H. Gonshor,
\emph{An Introduction to the Theory of Surreal Numbers},
London Mathematical Society Lecture Note Series 110,
Cambridge University Press, 1986.
\href{https://www.cambridge.org/core/books/an-introduction-to-the-theory-of-surreal-numbers/312AE504A3E88E804054BFB390446374}{Publisher record}.

\bibitem{Neumann}
B. H. Neumann,
On ordered division rings,
\emph{Transactions of the American Mathematical Society} \textbf{66}
(1949), 202--252.
\href{https://doi.org/10.1090/S0002-9947-1949-0032593-5}{doi:10.1090/S0002-9947-1949-0032593-5}.

\bibitem{BHW}
V. Benci, L. Horsten, and S. Wenmackers,
Non-Archimedean probability,
\emph{Milan Journal of Mathematics} \textbf{81} (2013), 121--151.
\href{https://doi.org/10.1007/s00032-012-0191-x}{doi:10.1007/s00032-012-0191-x}.
Preprint \href{https://arxiv.org/abs/1106.1524}{arXiv:1106.1524};
preprint Section 3.1 was used for the axiom comparison.

\bibitem{LK}
S. Ludkovsky and A. Khrennikov,
\emph{Stochastic processes on non-Archimedean spaces with values in
non-Archimedean fields}, 2001.
\href{https://arxiv.org/abs/math/0110305}{arXiv:math/0110305}.
Definition 2.2 and Theorem 2.15.2 in the retrieved version.

\bibitem{Repo}
V. Reshetnikov, \emph{Surreal}, GitHub repository,
snapshot \texttt{4cf691c7d951}, accessed 22 September 2026.
\href{https://github.com/VladimirReshetnikov/Surreal/tree/4cf691c7d951e037739d32d9f5c387dcce724f3c}{Full pinned snapshot}.
See \cref{subsec:comparison,sec:audit} for the audit scope.

\bibitem{RepoPhysics}
\emph{Surreal scalars and spacetime}, repository research report,
\href{https://github.com/VladimirReshetnikov/Surreal/tree/4cf691c7d951e037739d32d9f5c387dcce724f3c/docs/physics/surreal-scalars-and-spacetime}{cataloged directory};
compared with the available prior text \texttt{surreal\_physics.pdf},
Section 13.3. The equal-infinitesimal-weight observation is prior background,
not a novelty claim of this manuscript.
\end{thebibliography}
}

\end{document}
